Las tres boletas suman 100 puntos cada una, para un total de 300. El orden se construye de mayor a menor total, usando rango competitivo con empate: los requisitos con igual total comparten el mismo rango y el siguiente salta tantas posiciones como elementos empatados.

\begin{table}[htbp]
\centering
\footnotesize
\caption{Orden de prioridad según el total de las tres boletas.}
\label{tab:ranking_votacion}
\begin{tabular}{@{}ll@{}}
\toprule
Rango (total) & Requisitos \\
\midrule
1 (24) & R07 \\
2 (21) & R01, R10, R14, R17 \\
6 (18) & R03, R06, R19 \\
9 (15) & R02, R12 \\
11 (14) & R18 \\
12 (13) & R04 \\
13 (12) & R09, R13, R15 \\
16 (9) & R05, R08, R11, R16, R20 \\
\bottomrule
\end{tabular}
\end{table}

R04 (Must) obtuvo 13 puntos, el más bajo de su grupo; R03 (18) y R18 (14), ambos Could, igualan o superan a varios Must. R20 (Won't have) quedó en el tramo más bajo.

